\subsection{Sprint 1}

Após a entrega dos protótipos de telas para a cliente, nossos AGES IV montaram um quadro \textit{Kanban} no qual iria conter as US e tasks a serem trabalhadas ao longo do semestre. Além disso, o time foi dividido em três squads: duas focadas em desenvolvimento das US que deveriam ser entregues, e uma dedicada  a testes e revisões de \ac{mr}. Eu acabei sendo alocado para squad de testes.

Fiquei encarregado inicialmente, junto com meu colega AGES II Gabriel Suterio, de implementar o esquema do banco de dados usando o PrismaORM em nossa API e de gerar alguns dados iniciais de teste. Como nossos colegas AGES I da parte de backend estavam um pouco perdidos e ainda em processo de aprendizado, decidimos incorporá-los ao nosso esforço e trabalhar em grupo. 

O esquema do banco passou por algumas pequenas alterações ao longo da \textit{sprint}, incluindo ajustes na nomenclatura, tradução e tipos dos atributos.  Porém, logo que pronto, começamos definitivamente o desenvolvimento da API. Com o nosso mesmo grupo da tarefa anterior, trabalhamos na rota que retorna todos os eventos cadastrados dentro do aplicativo.
\\

\fbox{%
    \parbox{0.95\linewidth}{%
        \setlength{\parindent}{15pt}
            \itshape Quando o time foi dividido em squads e eu fui alocado na squad de testes, fiquei um pouco receoso. O escopo do projeto parecia menos complexo na parte de backend do que eu esperava, o que me deixou preocupado com a quantidade de trabalho a ser realizado. 

            Após esclarecer minha preocupação com os AGES IV, fui lembrado de que, quanto mais eu pudesse contribuir com outras tarefas, melhor seria para o projeto. Isso me deixou mais tranquilo e disposto a trabalhar no que fosse necessário para ajudar a equipe.

            A Sprint 1 foi relativamente tranquila, com exceção de alguns débitos técnicos. Esses problemas surgiram devido à alta demanda enfrentada pela equipe de frontend e à falta de materiais fornecidos pela cliente para tarefas específicas do backend.
    }%
}